% Options for packages loaded elsewhere
\PassOptionsToPackage{unicode}{hyperref}
\PassOptionsToPackage{hyphens}{url}
\documentclass[
]{article}
\usepackage{xcolor}
\usepackage[margin=1in]{geometry}
\usepackage{amsmath,amssymb}
\setcounter{secnumdepth}{-\maxdimen} % remove section numbering
\usepackage{iftex}
\ifPDFTeX
  \usepackage[T1]{fontenc}
  \usepackage[utf8]{inputenc}
  \usepackage{textcomp} % provide euro and other symbols
\else % if luatex or xetex
  \usepackage{unicode-math} % this also loads fontspec
  \defaultfontfeatures{Scale=MatchLowercase}
  \defaultfontfeatures[\rmfamily]{Ligatures=TeX,Scale=1}
\fi
\usepackage{lmodern}
\ifPDFTeX\else
  % xetex/luatex font selection
\fi
% Use upquote if available, for straight quotes in verbatim environments
\IfFileExists{upquote.sty}{\usepackage{upquote}}{}
\IfFileExists{microtype.sty}{% use microtype if available
  \usepackage[]{microtype}
  \UseMicrotypeSet[protrusion]{basicmath} % disable protrusion for tt fonts
}{}
\makeatletter
\@ifundefined{KOMAClassName}{% if non-KOMA class
  \IfFileExists{parskip.sty}{%
    \usepackage{parskip}
  }{% else
    \setlength{\parindent}{0pt}
    \setlength{\parskip}{6pt plus 2pt minus 1pt}}
}{% if KOMA class
  \KOMAoptions{parskip=half}}
\makeatother
\usepackage{color}
\usepackage{fancyvrb}
\newcommand{\VerbBar}{|}
\newcommand{\VERB}{\Verb[commandchars=\\\{\}]}
\DefineVerbatimEnvironment{Highlighting}{Verbatim}{commandchars=\\\{\}}
% Add ',fontsize=\small' for more characters per line
\usepackage{framed}
\definecolor{shadecolor}{RGB}{248,248,248}
\newenvironment{Shaded}{\begin{snugshade}}{\end{snugshade}}
\newcommand{\AlertTok}[1]{\textcolor[rgb]{0.94,0.16,0.16}{#1}}
\newcommand{\AnnotationTok}[1]{\textcolor[rgb]{0.56,0.35,0.01}{\textbf{\textit{#1}}}}
\newcommand{\AttributeTok}[1]{\textcolor[rgb]{0.13,0.29,0.53}{#1}}
\newcommand{\BaseNTok}[1]{\textcolor[rgb]{0.00,0.00,0.81}{#1}}
\newcommand{\BuiltInTok}[1]{#1}
\newcommand{\CharTok}[1]{\textcolor[rgb]{0.31,0.60,0.02}{#1}}
\newcommand{\CommentTok}[1]{\textcolor[rgb]{0.56,0.35,0.01}{\textit{#1}}}
\newcommand{\CommentVarTok}[1]{\textcolor[rgb]{0.56,0.35,0.01}{\textbf{\textit{#1}}}}
\newcommand{\ConstantTok}[1]{\textcolor[rgb]{0.56,0.35,0.01}{#1}}
\newcommand{\ControlFlowTok}[1]{\textcolor[rgb]{0.13,0.29,0.53}{\textbf{#1}}}
\newcommand{\DataTypeTok}[1]{\textcolor[rgb]{0.13,0.29,0.53}{#1}}
\newcommand{\DecValTok}[1]{\textcolor[rgb]{0.00,0.00,0.81}{#1}}
\newcommand{\DocumentationTok}[1]{\textcolor[rgb]{0.56,0.35,0.01}{\textbf{\textit{#1}}}}
\newcommand{\ErrorTok}[1]{\textcolor[rgb]{0.64,0.00,0.00}{\textbf{#1}}}
\newcommand{\ExtensionTok}[1]{#1}
\newcommand{\FloatTok}[1]{\textcolor[rgb]{0.00,0.00,0.81}{#1}}
\newcommand{\FunctionTok}[1]{\textcolor[rgb]{0.13,0.29,0.53}{\textbf{#1}}}
\newcommand{\ImportTok}[1]{#1}
\newcommand{\InformationTok}[1]{\textcolor[rgb]{0.56,0.35,0.01}{\textbf{\textit{#1}}}}
\newcommand{\KeywordTok}[1]{\textcolor[rgb]{0.13,0.29,0.53}{\textbf{#1}}}
\newcommand{\NormalTok}[1]{#1}
\newcommand{\OperatorTok}[1]{\textcolor[rgb]{0.81,0.36,0.00}{\textbf{#1}}}
\newcommand{\OtherTok}[1]{\textcolor[rgb]{0.56,0.35,0.01}{#1}}
\newcommand{\PreprocessorTok}[1]{\textcolor[rgb]{0.56,0.35,0.01}{\textit{#1}}}
\newcommand{\RegionMarkerTok}[1]{#1}
\newcommand{\SpecialCharTok}[1]{\textcolor[rgb]{0.81,0.36,0.00}{\textbf{#1}}}
\newcommand{\SpecialStringTok}[1]{\textcolor[rgb]{0.31,0.60,0.02}{#1}}
\newcommand{\StringTok}[1]{\textcolor[rgb]{0.31,0.60,0.02}{#1}}
\newcommand{\VariableTok}[1]{\textcolor[rgb]{0.00,0.00,0.00}{#1}}
\newcommand{\VerbatimStringTok}[1]{\textcolor[rgb]{0.31,0.60,0.02}{#1}}
\newcommand{\WarningTok}[1]{\textcolor[rgb]{0.56,0.35,0.01}{\textbf{\textit{#1}}}}
\usepackage{graphicx}
\makeatletter
\newsavebox\pandoc@box
\newcommand*\pandocbounded[1]{% scales image to fit in text height/width
  \sbox\pandoc@box{#1}%
  \Gscale@div\@tempa{\textheight}{\dimexpr\ht\pandoc@box+\dp\pandoc@box\relax}%
  \Gscale@div\@tempb{\linewidth}{\wd\pandoc@box}%
  \ifdim\@tempb\p@<\@tempa\p@\let\@tempa\@tempb\fi% select the smaller of both
  \ifdim\@tempa\p@<\p@\scalebox{\@tempa}{\usebox\pandoc@box}%
  \else\usebox{\pandoc@box}%
  \fi%
}
% Set default figure placement to htbp
\def\fps@figure{htbp}
\makeatother
\setlength{\emergencystretch}{3em} % prevent overfull lines
\providecommand{\tightlist}{%
  \setlength{\itemsep}{0pt}\setlength{\parskip}{0pt}}
\usepackage{bookmark}
\IfFileExists{xurl.sty}{\usepackage{xurl}}{} % add URL line breaks if available
\urlstyle{same}
\hypersetup{
  pdftitle={GroupH\_HM0},
  pdfauthor={A. Rossi, B. Verdi, C. Bianchi},
  hidelinks,
  pdfcreator={LaTeX via pandoc}}

\title{GroupH\_HM0}
\author{A. Rossi, B. Verdi, C. Bianchi}
\date{2022-10-20}

\begin{document}
\maketitle

{
\setcounter{tocdepth}{3}
\tableofcontents
}
\subsection{FSDS - Chapter 2}\label{fsds---chapter-2}

\subsubsection{Ex X.x}\label{ex-x.x}

\emph{Here the text of the first exercise.}

\textbf{Solution}

Add comments to the solution.

\begin{Shaded}
\begin{Highlighting}[]
\FunctionTok{summary}\NormalTok{(cars)}
\end{Highlighting}
\end{Shaded}

\begin{verbatim}
##      speed           dist       
##  Min.   : 4.0   Min.   :  2.00  
##  1st Qu.:12.0   1st Qu.: 26.00  
##  Median :15.0   Median : 36.00  
##  Mean   :15.4   Mean   : 42.98  
##  3rd Qu.:19.0   3rd Qu.: 56.00  
##  Max.   :25.0   Max.   :120.00
\end{verbatim}

\subsubsection{Ex X.x}\label{ex-x.x-1}

\emph{Here the text of the second exercise.}

\textbf{Solution}

Add comments to the solution.

\[ Y \sim N(0,1)\]

\subsection{CS - Chapter 1}\label{cs---chapter-1}

\subsubsection{Ex 1.1}\label{ex-1.1}

\emph{Here the text of the first exercise from Wood's textbook.}

\textbf{Solution}

Add comments to the solution.

\[ X \sim N(0,1)\]

\subsection{CS - Chapter 1}\label{cs---chapter-1-1}

\subsubsection{Ex 1.1}\label{ex-1.1-1}

\emph{Exponential random variable, \(X \ge 0\), has p.d.f.
\(f(x) = \lambda e^{-\lambda x}\).}

\begin{enumerate}
\def\labelenumi{\arabic{enumi}.}
\tightlist
\item
  \textbf{Find the c.d.f. and the quantile function for \(X\).}
\item
  \textbf{Find \(\Pr(X \le a)\) and the median of \(X\).}
\item
  \textbf{Find the mean and variance of \(X\).}
\end{enumerate}

\textbf{Solution}

\begin{enumerate}
\def\labelenumi{\arabic{enumi}.}
\item
  \textbf{Funzione di ripartizione (c.d.f.) e Quantile:} La funzione di
  ripartizione \(F(x)\) si ottiene integrando la densità di probabilità
  da \(0\) a \(x\):
  \[F(x) = P(X \le x) = \int_{0}^{x} \lambda e^{-\lambda t} dt = \left[ -e^{-\lambda t} \right]_{0}^{x} = 1 - e^{-\lambda x} \quad (\text{per } x \ge 0)\]

  Per trovare la funzione quantile \(F^{-1}(p)\), uguagliamo
  \(F(x) a p\) e risolviamo rispetto a \(x\):
  \[1 - e^{-\lambda x} = p \implies e^{-\lambda x} = 1 - p \implies -\lambda x = \ln(1-p)\]
  Da cui si ottiene la funzione quantile:
  \[Q(p) = -\frac{\ln(1-p)}{\lambda} \quad \text{per } p \in (0,1)\]
\item
  \textbf{Probabilità \(\Pr(X \le a)\) e Mediana:}

  \begin{itemize}
  \tightlist
  \item
    La probabilità che \(X\) sia minore o uguale a una costante
    \(a \ge 0\) è data direttamente dalla c.d.f.:
    \[\Pr(X \le a) = F(a) = 1 - e^{-\lambda a}\]
  \item
    La mediana corrisponde al quantile di ordine \(p = 0.5\):
    \[\text{Mediana} = Q(0.5) = -\frac{\ln(1 - 0.5)}{\lambda} = -\frac{\ln(0.5)}{\lambda} = \frac{\ln(2)}{\lambda}\]
  \end{itemize}
\item
  \textbf{Media e Varianza di \(X\):}

  \begin{itemize}
  \tightlist
  \item
    \textbf{Media \(E[X]\):} Integrando per parti:
    \[E[X] = \int_{0}^{\infty} x \lambda e^{-\lambda x} dx = \left[ -x e^{-\lambda x} \right]_{0}^{\infty} + \int_{0}^{\infty} e^{-\lambda x} dx = 0 + \left[ -\frac{e^{-\lambda x}}{\lambda} \right]_{0}^{\infty} = \frac{1}{\lambda}\]
  \item
    \textbf{Varianza \(\text{Var}(X)\):} Calcoliamo prima il momento
    secondo \(E[X^2]\) (sempre tramite integrazione per parti):
    \[E[X^2] = \int_{0}^{\infty} x^2 \lambda e^{-\lambda x} dx = \frac{2}{\lambda^2}\]
    Utilizzando la relazione fondamentale della varianza
    \(\text{Var}(X) = E[X^2] - (E[X])^2\):
    \[\text{Var}(X) = \frac{2}{\lambda^2} - \left(\frac{1}{\lambda}\right)^2 = \frac{2 - 1}{\lambda^2} = \frac{1}{\lambda^2}\]
  \end{itemize}
\end{enumerate}

\begin{center}\rule{0.5\linewidth}{0.5pt}\end{center}

\subsubsection{Ex 1.6}\label{ex-1.6}

\emph{Let \(X\) and \(Y\) be random variables with
\(\text{Var}(X) = \sigma_x^2\), \(\text{Var}(Y) = \sigma_y^2\), and
\(\text{Cov}(X,Y) = \sigma_{xy}\). Find \(\text{Var}(X+Y)\) and
\(\text{Var}(X-Y)\).}

\textbf{Solution}

Ricordando le proprietà della varianza per la somma di due variabili
casuali (non necessariamente indipendenti):

\begin{enumerate}
\def\labelenumi{\arabic{enumi}.}
\item
  \textbf{Varianza della somma \(\text{Var}(X + Y)\):}
  \[\text{Var}(X + Y) = \text{Var}(X) + \text{Var}(Y) + 2\text{Cov}(X,Y) = \sigma_x^2 + \sigma_y^2 + 2\sigma_{xy}\]
\item
  \textbf{Varianza della differenza \(\text{Var}(X - Y)\):}
  \[\text{Var}(X - Y) = \text{Var}(X) + \text{Var}(-Y) + 2\text{Cov}(X, -Y) = \sigma_x^2 + \sigma_y^2 - 2\sigma_{xy}\]
  \emph{(poiché \(\text{Var}(-Y) = (-1)^2\sigma_y^2 = \sigma_y^2\) e
  \(\text{Cov}(X, -Y) = -\text{Cov}(X,Y)\))}
\end{enumerate}

\end{document}
